\documentclass[12pt, letterpaper]{article}
\usepackage[utf8]{inputenc}
\usepackage[spanish]{babel}
\usepackage{amsmath, amssymb, amsthm, bm}
\usepackage{graphicx}
\usepackage{array}
\usepackage{enumitem}
\usepackage{float}
\usepackage{listings}
\usepackage{xcolor}
\usepackage{tabularx}
\usepackage{booktabs}
%% Sets page size and margins
\usepackage[top=2.5cm,bottom=2.5cm,left=2cm,right=2cm]{geometry}
\usepackage{setspace}
\onehalfspacing
\usepackage{gensymb}
\usepackage{fancyhdr}
\pagestyle{fancy}
\fancyhead[L]{Compiladores}
\fancyhead[C]{Práctica 4}
\fancyhead[R]{Facultad de Ciencias UNAM}
\fancyfoot[C]{\thepage}
\renewcommand{\headrulewidth}{0.4pt}
\renewcommand{\footrulewidth}{0.4pt}

\lstdefinestyle{bashstyle}{
    backgroundcolor=\color{gray!10},
    basicstyle=\small\ttfamily,
    breaklines=true,
    frame=single,
    frameround=tttt
}

\lstset{
    inputencoding=utf8,
    extendedchars=true,
    literate=
        {á}{{\'a}}1
        {é}{{\'e}}1
        {í}{{\'i}}1
        {ó}{{\'o}}1
        {ú}{{\'u}}1
        {Á}{{\'A}}1
        {É}{{\'E}}1
        {Í}{{\'I}}1
        {Ó}{{\'O}}1
        {Ú}{{\'U}}1
        {ñ}{{\~n}}1
        {Ñ}{{\~N}}1
        {ü}{{\"u}}1
        {Ü}{{\"U}}1
}


\title{
    \vspace*{-2.5cm}
    \setlength{\parindent}{0cm} %Indentation 0
    \rule{\linewidth}{0.1mm}
    \vspace*{-1.4cm}
    \begin{center}
        \begin{minipage}{2.5cm}
            \begin{center}
            \includegraphics[height=3.4cm]{images/Logo_UNAM.png}
            \end{center}
        \end{minipage}\hfill
        %--------------
        \hspace{0.3cm}
        \begin{minipage}{11cm}
            \begin{center}
            \normalfont \small \textsc{UNIVERSIDAD NACIONAL AUTÓNOMA DE MÉXICO \\ FACULTAD DE CIENCIAS \\ Compiladores } \\
    		\normalsize \textbf{Práctica 4}\\
            \normalsize \textsc{Flex}\\
            \end{center}
        \end{minipage}\hfill
        %--------------
        \begin{minipage}{3cm}
            \begin{center}
            \includegraphics[height=3.6cm]{images/Logo_FC.png}
            \end{center}
        \end{minipage}
    \end{center}
    \vspace*{-0.25cm}
    \rule{\linewidth}{0.1mm}
    \vspace*{0.5cm}
    {\normalsize \textbf{Integrantes:} \par}
    \vspace{-0.5cm}
    \begin{tabular}{{l@{\hspace{1cm}}r}}
        \vspace*{-0.5cm}
        \small Escobar Gonzalez Isaac Giovani & \small 321336400 \\
        \vspace*{-0.5cm}
        \small Garduño Escobar Kevin Jonathan & \small 321070629 \\
        \vspace*{-0.5cm}
        \small González Gutiérrez Antonio Tonatiuh & \small 321195476 \\
        \vspace*{-0.5cm}
        \small Zaldivar Alanis Rodrigo & \small 424029605
    \end{tabular}
    \date{}
}

\begin{document}
\maketitle
\vspace*{-2.5cm}
\subsection*{Estrategia de diseño léxico y condiciones de inicio}

\subsection*{Instrucciones de compilación y ejecución de pruebas}

\subsection*{Bitácora de casos límite y depuración léxica}

Para la parte de Operadores y Delimitadores, el principal caso límite enfrentado fue la resolución de ambigüedades en secuencias consecutivas de símbolos matemáticos, como el lexema \texttt{+++} o la diferenciación entre \verb|>| y \verb|>=|. Este problema radicaba en que Flex podía interpretar anticipadamente caracteres individuales y dividir operadores compuestos si las reglas no estaban estructuradas lógicamente. La solución implementada consistió en aprovechar la regla de la coincidencia más larga (\textit{longest match}) natural de Flex, estructurando el bloque de operadores en \texttt{scanner.l} en un orden estricto de longitud descendente. Se colocaron primero las asignaciones compuestas de 3 caracteres (como \verb|<<=|), seguidas por las de 2 caracteres (como \texttt{++}, \texttt{+=}) y, finalmente, los caracteres simples (\texttt{+}, \texttt{=}). Esto garantizó que lexemas de prueba como \texttt{+=+} se reconocieran correctamente como \texttt{TOK\_PLUS\_ASSIGN} seguido de \texttt{TOK\_PLUS}.

\subsection*{Limitaciones conocidas y deuda técnica}

En cuanto a los operadores a nivel de bits, una limitación del diseño léxico es que se modelaron como palabras de texto (\texttt{convergencia}, \texttt{expansion}, \texttt{oscuridad}) para adaptarse a la temática del lenguaje \textit{CosmoLang}. Debido a que Flex evalúa mediante expresiones regulares sin contexto sintáctico, estas palabras coinciden idénticamente con el patrón definido para los identificadores (\texttt{ID}). Aunque el conflicto se resolvió declarando las reglas de los operadores literales en bloques superiores a la regla de identificadores, la deuda técnica podría ser que los usuarios del lenguaje tienen estrictamente prohibido utilizar palabras como \texttt{expansion} o \texttt{superposicion} para nombrar variables o funciones, ya que el escáner las interceptará forzosamente como operadores, limitando el vocabulario disponible para el programador.

\newpage
\subsection*{Validación de Trabajo Presencial / Reposición Teórica}
No se que va a aquí xd

\begin{itemize}
    \item \textbf{Isaac Giovani Escobar González}:
    \item \textbf{Kevin Jonathan Garduño Escobar}:
    \item \textbf{Antonio Tonatiuh González Gutiérrez}:
    \item \textbf{Rodrigo Zaldivar Alanis}:
\end{itemize}

\newpage
\subsection*{Evaluación de la Práctica}
\noindent Grado de acuerdo con cada una de las siguientes afirmaciones utilizando una escala de 1 a 5, donde:
\begin{itemize}
    \item 1: Totalmente en desacuerdo
    \item 2: En desacuerdo
    \item 3: Neutral
    \item 4: De acuerdo
    \item 5: Totalmente de acuerdo
\end{itemize}

\begin{table}[H]
    \centering
    \begin{tabularx}{\textwidth}{|X|c|c|c|c|c|}
        \hline
        \textbf{Criterio / Afirmación} & \textbf{1} & \textbf{2} & \textbf{3} & \textbf{4} & \textbf{5} \\
        \hline
        Los conceptos y habilidades desarrollados en esta práctica son relevantes para mi formación profesional. &  &  & \checkmark &  &  \\
        \hline
        El trabajo realizado me permitió comprender de forma práctica la teoría revisada en clase. &  &  &  & \checkmark &  \\
        \hline
        Las actividades y problemáticas planteadas durante la sesión resultaron atractivas y motivadoras. &  &  & \checkmark &  &  \\
        \hline
        El enfoque práctico y los retos propuestos captaron mi interés a lo largo del desarrollo. &  &  &  & \checkmark &  \\
        \hline
        El nivel de complejidad técnica fue adecuado respecto a los conocimientos y recursos previos disponibles. &  &  &  & \checkmark &  \\
        \hline
        La claridad de la guía permitió realizar la práctica de manera fluida y sin bloques innecesarios. &  &  & \checkmark &  &  \\
        \hline
    \end{tabularx}
\end{table}

\end{document}
